\thispagestyle{plain}

\pdfbookmark[1]{Abstract}{abstract}
\chapter*{Abstract}
This project explores the implementation of digital communication systems using the software radio technique. By shifting traditional hardware-based radio problems into the software domain, we simulate a transmitter and receiver environment that closely mirrors real-world antenna behavior. The primary focus is the generation and analysis of three specific line codes: Polar RZ, Polar NRZ, and Unipolar.

Since the transmitted bit-streams and initial timing offsets are inherently
random, the system is modeled as a random process. Using MATLAB, we generate an
ensemble of realizations for each signaling technique to evaluate their
statistical and time-domain characteristics. Specifically, we analyze the
statistical mean, time mean, autocorrelation functions, and power spectral
density. These results are then used to verify the stationarity and ergodicity
of the processes, providing a comprehensive understanding of how
software-defined waveforms behave in a practical communications environment.